% 加入 x1<=2（松弛x5）后，可行域不变，但顶点(2,0)上有 x2=0、2x1+x2=4、x1=2 三条边界同时取等号。
% 基{1,4,5}、{1,3,4}、{1,2,4}都表示这一点；从{1,3,4}入基x2时最小比值为0，是退化换基；随后入基x5到(4/3,4/3)。
\begin{tikzpicture}[x=14mm,y=14mm,
  note/.style={font=\FigureNoteFont\fontsize{7.5}{9}\selectfont,text=FigureMuted}]
  \fill[FigureBlueFill] (0,0)--(2,0)--(4/3,4/3)--(0,2)--cycle;
  \draw[FigureBlue!60,line width=.6pt] (0,0)--(2,0)--(4/3,4/3)--(0,2)--cycle;
  \draw[FigureStroke!70,line width=.6pt] (1.25,1.5)--(2.25,-0.5);
  \draw[FigureStroke!70,line width=.6pt] (2,-0.5)--(2,1.6);
  \draw[->,FigureStroke,line width=.6pt] (-0.15,0)--(2.8,0) node[right]{$x_1$};
  \draw[->,FigureStroke,line width=.6pt] (0,-0.15)--(0,2.3) node[above]{$x_2$};
  \draw[-{Stealth[length=2.4mm]},FigureRed,line width=1.2pt] (0.05,0)--(1.9,0);
  \draw[-{Stealth[length=2.2mm]},FigureRed,line width=1pt] (2.07,0.07) arc[start angle=-60,end angle=250,radius=0.14];
  \draw[-{Stealth[length=2.4mm]},FigureRed,line width=1.2pt] (1.97,0.06)--(1.38,1.26);
  \fill[FigureRed] (2,0) circle (1.8pt);
  \fill[FigureInk] (4/3,4/3) circle (1.5pt);
  \fill[FigureInk] (0,0) circle (1.5pt);
  \node[note,anchor=south] at (2.1,1.6) {$x_1=2$（$x_5=0$）};
  \node[note,anchor=east] at (1.22,1.62) {$2x_1+x_2=4$（$x_3=0$）};
  \node[note,anchor=north west,align=left] at (2.08,-0.08) {同一点$(2,0)$的三个基：\\$\{1,4,5\}$，$\{1,3,4\}$，$\{1,2,4\}$};
  \node[note,anchor=south west,text=FigureRed!80!FigureInk] at (2.2,0.3) {零步长换基};
  \draw[FigureMuted,line width=.4pt] (2.22,0.38)--(2.12,0.25);
  \node[note,anchor=north east] at (1.28,1.3) {$(4/3,4/3)$};
\end{tikzpicture}
